%===================================================================================================
\section{Observed Issues in Existing Compiler/Test Output} %++++++++++++++++++++++++++++++++++++++++
  \label{Observed Issues in Existing Compiler/Test Output} %++++++++++++++++++++++++++++++++++++++++
%===================================================================================================


\improve


\begin{justify}
This section identifies recurring issues in baseline outputs produced by common compiler and testing
tools. The focus is on how test failures are reported and how students must reconstruct key
information that is not explicitly presented. The examples are drawn from representative beginner
exercises used throughout this thesis.
\end{justify}


%===============================================================================
\subsection{Examples of current error messages} %+++++++++++++++++++++++++++++++
     \label{Examples of current error messages} %+++++++++++++++++++++++++++++++
%===============================================================================


\begin{justify}
Error messages often require the programmer to understand not only \emph{what} was tested but also
\emph{how} it was tested. For instance, tests for data structures and stateful components may
involve sequences of operations (e.g., insertions and removals) followed by an invariant check.
If a failure occurs after multiple steps, baseline output rarely describes the sequence leading to
the assertion point, even though a compact description of the minimal failing history would often be
sufficient for debugging.

In addition, baseline frameworks frequently present essential information in a fragmented manner.
For example, \texttt{Assert.AreEqual} can report an expected/actual mismatch, but it does not
inherently encode the tested expression. Unless the test method name and surrounding context are
highly descriptive, students may struggle to associate the mismatch with a concrete call in their
own code. As a consequence, learners must infer missing context from indirect cues such as test
names, inputs printed elsewhere, or stack traces.
\end{justify}


%===========================================================
\subsubsection{Example 1 --- \texttt{digitSum}} %+++++++++++
        \label{Example 1 digitSum} %++++++++++++++++++++++++
%===========================================================


\begin{justify}
Consider the exercise \fsharp{digitSum : Nat |$\to$| Nat}, where the task is to compute the digit
sum of a natural number. A representative novice mistake is an incorrect base case, such as
returning \texttt{1N} for input \texttt{0N}. Example-based tests using Unquote can display the
failed equality as a quoted expression and show the boolean result of the comparison, but the output
also contains repeated, qualified expression renderings and a stack trace that dominates the overall
message.

For novices, this reduces the signal-to-noise ratio: the crucial mismatch is present, yet not
highlighted as a compact summary.

When property-based testing is used (e.g., FsCheck integrated with MSTest assertions), the output
typically reports that a property is falsifiable and provides an original and a shrunk
counterexample (often a minimal input such as \texttt{0N}). This information is valuable, but the
message does not necessarily provide a direct, novice-friendly statement of the failing behaviour in
terms of a single reconstructed expression such as:

\begin{quote}
    \emph{Tested expression: \texttt{digitSum 0N}. Expected: \texttt{0N}. Actual: \texttt{1N}.}
\end{quote}

Instead, students must combine fragments from different parts of the output (counterexample,
expected/actual mismatch, and implicit knowledge about what the test invokes) to arrive at that
statement.
\end{justify}


%===========================================================
\subsubsection{Example 2 --- \texttt{sortedDigits}} %+++++++
        \label{Example 2 sortedDigits} %++++++++++++++++++++
%===========================================================


\begin{justify}
A similar pattern occurs for boolean-valued exercises such as \\
\fsharp{sortedDigits : Nat |$\to$| bool}. Failures may hinge on subtle predicate choices
(e.g., strict vs.\ non-strict comparisons), where the difference between \texttt{<} and \texttt{<=}
becomes visible only on specific counterexamples (e.g., repeated digits). In baseline outputs, the
essential information is the concrete failing input together with the expected and actual boolean
result. However, this information is often presented without an explicit reconstruction of the
tested call, and it is not consistently formatted across different testing mechanisms \cite{barik}.
As a result, students must infer which invocation produced the mismatch and why the counterexample
is relevant to the specification.
\end{justify}


%===============================================================================
\subsection{Recurring structural problems} %++++++++++++++++++++++++++++++++++++
     \label{Recurring structural problems} %++++++++++++++++++++++++++++++++++++
%===============================================================================


Across the examples above, the following structural problems recur.


%===========================================================
\paragraph{Missing explicit linkage to the tested expression.}
    \label{Missing explicit linkage to the tested expression.}
Many assertion messages focus on the mismatch (expected vs.\ actual) without explicitly reporting
the expression or call that produced it. For novices, this disconnect is critical because it
introduces ambiguity about what exactly was evaluated. The problem is amplified when tests are
hidden, since the student cannot inspect the failing test case and must rely on the message alone.


%===========================================================
\paragraph{Fragmentation of essential information.}
    \label{Fragmentation of essential information.}
Even when all necessary facts are present (inputs, expected value, actual value), they may appear in
separate lines or in different parts of the output. Students are required to reconstruct the failing
behaviour by manually combining fragments. This reconstruction increases cognitive load and
encourages trial-and-error rather than systematic debugging.


%===========================================================
\paragraph{Inconsistent message formats across tools.}
    \label{Inconsistent message formats across tools.}
In the considered setting, example-based tests (e.g., Unquote) and framework assertions (e.g.,
MSTest \texttt{Assert.AreEqual}) produce different output structures and emphasize different aspects
of a failure. Such inconsistency prevents novices from developing a stable reading strategy.
More generally, diagnostics often originate from internal tool representations that do not map
cleanly back to what the programmer wrote \cite{becker}.


%===========================================================
\paragraph{Verbosity and low-level noise.}
    \label{Verbosity and low-level noise.}
Stack traces, framework invocation details, and module-qualified names can obscure the central
mismatch. While such details may help experienced developers diagnose infrastructure issues, they
often do not help novices fix solution logic. Verbose output can reduce readability and make it
difficult to locate actionable information.


%===========================================================
\paragraph{Locality and causality gaps.}
    \label{Locality and causality gaps.}
Error reports are not always local to the underlying cause. Becker et al.\ note that errors are
frequently detected far away from their generation sites (the ``locality problem''), and that
diagnostics may reflect tool-internal viewpoints rather than the programmer's conceptual model
\cite{becker}. In educational contexts, this manifests when the message points to the test harness
location (which students may not access) rather than to the relevant part of their own code.

This section characterises the structural properties of baseline outputs that contribute to poor
interpretability for novices. The analysis is organised around recurring patterns, illustrated
through representative examples from beginner exercises.


%===============================================================================
\subsection{Issue 1: Essential information is fragmented or implicit} %+++++++++
     \label{Issue 1: Essential information is fragmented or implicit} %+++++++++
%===============================================================================


\begin{justify}
In many baseline test outputs, the information required to reconstruct the failing behaviour is not
presented as a single, explicit summary. Instead, learners must infer it by combining multiple
fragments (e.g., a test name, a partially printed expression, and a comparison line). When the
\emph{tested expression} is not explicitly shown, students may not be able to associate the failure
with a specific function call in their submission, especially when a test method executes multiple
assertions or multiple inputs.
\end{justify}


%===============================================================================
\subsection{Issue 2: Mixing core mismatch information with framework internals}%
     \label{Issue 2: Mixing core mismatch information with framework internals}%
%===============================================================================


\begin{justify}
	Failure reports frequently interleave the core mismatch (expected vs. actual) with verbose framework-specific content such as stack traces, internal runner frames,
	and reflection invocation details. While such information can be valuable for expert debugging, it often dominates the output visually and makes it difficult for novices to
	identify the actionable part of the report.
\end{justify}


%===============================================================================
\subsection{Issue 3: Inconsistent structure across test types and frameworks} %+
     \label{Issue 3: Inconsistent structure across test types and frameworks} %+
%===============================================================================


\begin{justify}
Even within a single course setup, students may encounter multiple output formats depending on
whether a failure stems from:

\skipS\begin{enumerate}[label={(\alph*)}]
    \item a unit-test assertion,
    \item a quotation-based assertion mechanism, or
    \item a property-based test with counterexample generation and shrinking.
\end{enumerate}

\skipS If these formats differ substantially in vocabulary, ordering, and level of detail, students
cannot form stable expectations about where to look for essential information.
\end{justify}


%===============================================================================
\subsection{Issue 4: Property-based testing output is informative but not pedagogically structured}
     \label{Issue 4: Property-based testing output is informative but not pedagogically structured}
%===============================================================================


\begin{justify}
Property-based testing frameworks can provide powerful information (e.g., a counterexample input and
a shrunk minimal failing case). However, this information is often embedded in output that assumes
familiarity with the property-testing workflow. In particular, the relationship between (i) the
generated input, (ii) the shrunk counterexample, and (iii) the failing assertion is not always made
explicit in a way that supports novice mental models.
\end{justify}


%===============================================================================
\subsection{Issue 5: Missing execution context for stateful or multi-step scenarios}
     \label{Issue 5: Missing execution context for stateful or multi-step scenarios}
%===============================================================================


\begin{justify}
For exercises involving stateful data structures or sequences of operations, a single
expected/actual mismatch may be insufficient to understand the failure. If the failure occurs after
a sequence of updates, students may need to know \emph{what was executed before the failing check}
in order to reproduce the bug. Baseline outputs typically do not preserve such operational context.

\skipM\textbf{Checkpoint (examples and evidence to add):}

\skipS\begin{itemize}
    \item TODO: Insert 2--3 concrete baseline outputs from your own exercise set and annotate them
        (highlight: where is the tested expression? where is expected/actual? what is noise?).
    \item TODO: Add at least one example for each: (i) example-based unit test failure, (ii)
        quotation/assertion failure, (iii) property-based failure with shrinking.
    \item TODO: Provide a short ``signal vs.\ noise'' breakdown for each example (e.g., mark lines
        as essential vs.\ auxiliary).
    \item TODO: If you discuss multi-step/stateful tests, add a small scenario (e.g., linked list
        operations) and show precisely what context is missing in baseline output.
\end{itemize}
\end{justify}


%===============================================================================
\subsection{Mock-up: additional observations to validate with literature (not yet backed by your local examples)}
     \label{Mock-up: additional observations to validate with literature (not yet backed by your local examples)}
%===============================================================================


\begin{justify}
The following patterns are frequently discussed in the broader literature on novice-facing
diagnostics and may apply to your setting, but should be validated with either (i) your own
collected outputs or (ii) clearly cited prior work:

\skipS\begin{itemize}
    \item TODO: \textbf{Jargon and modality switching:} natural-language fragments mixed with code
        tokens can force frequent context switches during reading.
    \item TODO: \textbf{Redundant repetition:} repeating the same information in multiple forms
        (e.g., multiple representations of the same expression) can increase reading effort without
        improving comprehension.
    \item TODO: \textbf{Poor localisation:} messages may be accurate but point novices to symptoms
        rather than the likely source of misunderstanding.
\end{itemize}
\end{justify}


%===============================================================================
% TODO +++++++++++++++++++++++++++++++++++++++++++++++++++++++++++++++++++++++++
%===============================================================================


% TODO


%===================================================================================================
% EOF ++++++++++++++++++++++++++++++++++++++++++++++++++++++++++++++++++++++++++++++++++++++++++++++
%===================================================================================================